% Section 2 — Background and Motivation
% Six subsections covering everything a reader needs before §3 (the agent):
%   2.1 The incident-triage workflow and where automation breaks
%   2.2 The data gap and why a controlled lab is required
%   2.3 Retrieval over institutional memory: the design space we draw from
%   2.4 Tool-using language-model agents and the open question of tool-use lift
%   2.5 Cost-aware diagnostic systems
%   2.6 Multi-window alert deduplication
% Every citation key resolves to a populated entry in references.bib
% (verified 2026-06-16). TAWOS issue-count claim cross-checked against
% DOCS/docs6/RESEARCH-CHARTER.md (458k issues; project count not
% documented in docs8 and therefore not asserted).

\section{Background and Motivation}
\label{sec:background}

\subsection{The incident-triage workflow}
\label{sec:background:workflow}

A production microservice cluster emits four kinds of evidence on every request: \emph{structured logs} from each service's request-handling and error paths; \emph{distributed traces} that record the per-span timing and error attributes of every dependency call; \emph{numeric metrics} that report request rates, error rates, and latency percentiles per service per second; and \emph{Kubernetes events} that capture pod lifecycle and infrastructure state changes. Together this telemetry is the substrate on which automated anomaly detection has matured: histogram baselines, learned classifiers on numeric features, and statistical alert rules answer \emph{``is something wrong?''} reliably enough that detection is no longer the binding constraint on response time~\cite{notaro2021aiops}.

What is binding is the step that follows. When an alert fires, the on-call engineer must decide three things: (i) is this a real incident or a benign spike; (ii) does it match a past incident the team has already diagnosed; and (iii) if so, which past Jira ticket describes it? The third question is the manual lookup we automate in this paper. In a mature organization the Jira corpus accumulates over years and represents the team's accreted understanding of how its system fails: each ticket records, in free text, the symptoms that prompted the page, the engineer's diagnosis, the underlying root cause, and the eventual fix~\cite{jira}. Because tickets are not indexed against telemetry, today's lookup is performed by keyword search; in practice the search takes ten to thirty minutes per page and depends on the engineer recalling the right vocabulary. A retrieval system that lifts the right past ticket from telemetry alone is therefore worth measuring, and is the operational target of this work.

\subsection{The data gap}
\label{sec:background:datagap}

Direct evaluation of telemetry-to-ticket retrieval is blocked by a data gap. The widely-used public log datasets---HDFS~\cite{xu2009hdfs} and the BlueGene/L and Thunderbird system logs~\cite{oliner2007supercomputers}---contain rich log streams but no ticket-side ground truth: there is no record of which past incident any given anomaly window corresponds to, and so they evaluate detection rather than diagnosis. Conversely, the largest public Jira corpus we are aware of contains roughly $458{,}000$ real issues from open-source projects but no co-collected telemetry: it is useful for ticket-side language modeling but not for joint retrieval. Real production telemetry from cooperating organizations is rarely shareable for legal and competitive reasons, is almost never reproducible by external reviewers, and is hard to share at the level of granularity (logs \emph{and} traces \emph{and} metrics \emph{and} Kubernetes events, all timestamped at sub-second resolution and joined to ticket identifiers) that this evaluation requires.

We therefore construct a controlled lab in which telemetry and Jira are generated jointly under known scenario labels. The lab forks a public microservice benchmark~\cite{onlineboutique}, instruments it with production-realistic OpenTelemetry SDKs across all services~\cite{otel-spec}, defines a taxonomy of incident scenario families, and executes the taxonomy by automated fault injection while a separate process emits Jira-style tickets describing each incident in engineer voice. To stress external validity we also build a second telemetry-rich corpus on a structurally different polyglot application~\cite{otel-demo}, and add a third corpus consisting of real Apache Jira tickets drawn from public bug trackers. The three corpora are described in detail in \S\ref{sec:evaluation:datasets}; what matters at this point is that the joint construction is reproducible end-to-end by any external reviewer with access to a Kubernetes cluster, and is generated under a strict label-leakage discipline so that the telemetry channel cannot reveal the scenario taxonomy to a model that should be learning it from the joint distribution.

\subsection{Retrieval over institutional memory}
\label{sec:background:retrieval}

The closest established research area is text-retrieval over an institutional memory corpus. Three families of retriever dominate the literature, and each one's strengths and weaknesses recur when the corpus is Jira-shaped.

\emph{Sparse retrieval} indexes documents by token overlap, weighted by inverse-document-frequency variants such as BM25~\cite{robertson2009bm25}. It is cheap, interpretable, and strong when query and document share vocabulary. It is weak when the same incident is described in different words by the engineer who wrote the ticket and the telemetry that surfaces a new instance.

\emph{Dense retrieval} encodes queries and documents into a common vector space via a pre-trained Transformer encoder, optionally fine-tuned on in-domain pairs~\cite{reimers2019sbert,karpukhin2020dpr}. It captures paraphrase and acronym-vs-expansion robustly, but its embeddings are opaque and its accuracy on a new domain depends on fine-tuning data. SPLADE~\cite{formal2021splade} is a sparse-but-learned variant that combines dense expressiveness with sparse index efficiency.

\emph{Hybrid retrieval} combines a sparse and a dense ranker (and optionally a knowledge-graph ranker) via late fusion. Reciprocal Rank Fusion~\cite{cormack2009rrf} is the canonical recipe: each ranker emits a per-document score, ranks are linearly combined under a fixed denominator, and ties are broken by per-ranker rank. Hybrid fusion is the de facto winner on many retrieval benchmarks because the failure modes of sparse and dense retrievers tend to be orthogonal, but the choice of fusion weights is application-specific.

\emph{Cross-encoder reranking} adds a second-stage Transformer that jointly attends to a query and a candidate, producing a single relevance score per pair~\cite{nogueira2019rerank}. Rerankers improve top-$k$ ordering meaningfully but cost an order of magnitude more compute per call than first-stage retrieval, which limits how deep the candidate pool can be in production.

A diagnostic system that aspires to robust real-world coverage has to draw on several of these primitives, because no single retriever dominates on every signal type: a numeric-feature anomaly looks different from a Kafka outage that surfaces in trace error spans, which looks different again from an authentication blip that surfaces in unordered log fragments inside a Jira ticket comment. The question this paper takes up is not how to build a better retriever in isolation, but how to compose existing retrievers \emph{per window} so that only the relevant primitives fire.

\subsection{Tool-using language-model agents}
\label{sec:background:agents}

Language models can be augmented at inference time with the ability to call external tools---retrievers, APIs, calculators, code interpreters---and to interleave tool-calls with reasoning steps. The \textsc{ReAct} framework~\cite{yao2023react} formalized this loop as an alternating sequence of \emph{reason} steps (free-form chain-of-thought) and \emph{act} steps (structured tool calls), terminating when the model emits a final answer. Subsequent work has explored which tools to expose, how to schedule calls, and how to learn the policy that decides which tool to call when~\cite{schick2023toolformer,lewis2020rag,izacard2023atlas}.

When the goal is diagnosis rather than answering a factual question, tool-using agents are conceptually well matched to the problem: a window where a numeric classifier is unsure can prompt the agent to fetch additional evidence (a wider trace window, the recent pod-event stream, an excerpt of past tickets from the same service) before committing to a top-$k$ ranking. The conceptual appeal is large enough that the term \emph{agentic retrieval} has begun to appear in the literature for exactly this pattern. Whether on-demand evidence gathering produces measurable retrieval-quality lift over a well-tuned non-agentic retriever---and at what cost---is an open empirical question; this paper measures it directly with paired-bootstrap confidence intervals across three datasets, and is candid about both the headline gains and the regimes in which lift is smaller than the prior literature suggests.

\subsection{Cost-aware diagnostic systems}
\label{sec:background:cost}

Inference cost has been studied as a first-class design variable since the cascaded face-detector literature, where Viola and Jones~\cite{viola2001cascade} showed that an asymmetric pipeline---cheap classifiers first, expensive classifiers only on cheap-positive examples---can reduce wall-time by orders of magnitude with negligible accuracy loss. The same idea recurs in deep-learning early-exit networks~\cite{teerapittayanon2016branchynet} and in retrieval-augmented systems that gate an expensive cross-encoder behind a cheap dense first-stage~\cite{nogueira2019rerank}.

In incident triage, cost-aware design has been articulated less explicitly. The dominant operational pattern is what we will call the \emph{always-on cascade}: every diagnostic pipeline (every retriever, every language-model verifier, every knowledge-graph query) is run on every window so that the most expensive signal is available if needed. This pattern is robust, simple to deploy, and rapidly expensive: a language-model verifier that costs $\$0.0005$ per call at gpt-4o-mini rates is $\$0.50$ per thousand windows even before token-budget overruns. The alternative---a single fixed pipeline that omits the expensive verifier---trades robustness for cost. Neither pattern adapts to the evidence each window actually carries. The cost-saving opportunity that this paper measures is the gap between the always-on cascade and a per-window adaptive policy that invokes the verifier only when the cheap path is uncertain.

\subsection{Multi-window alert deduplication}
\label{sec:background:suppression}

A real incident does not occupy a single monitoring window. A cascading database degradation may produce alerts on consecutive five-minute windows for an hour before the system recovers; a flapping pod may emit alerts intermittently for days. Treating each window as an independent paging decision produces alert fatigue, a well-documented operational hazard that erodes trust in monitoring and increases the time-to-acknowledge of new incidents~\cite{notaro2021aiops}. Engineering practice deduplicates by some combination of service identity, scenario family, and recent monitoring state, but the published literature on the exact suppression policy is thin: the rules tend to live in runbooks rather than papers. We adopt a conservative rule (\S\ref{sec:approach:state}) that suppresses a page when the same top-ranked ticket has been emitted on each of the last three contiguous windows for the same service and scenario family, provided no recovery window has intervened. The operational metric we use to evaluate this rule, \emph{pages per incident}, is the number of paged tickets divided by the number of distinct incident identifiers and is the metric on-call engineers report caring about most.
